%---------------------------------------------------------------
\chapter{Rešerše}
\label{ch:reserse}
%---------------------------------------------------------------

\section{Biomechanika plavání}
\label{sec:biomechanika}
% TODO: Styly, metriky, co dělá dobrého plavce
% Zdroje: Maglischo 2003, Craig & Pendergast 1979, Barbosa 2011
% - 4 plavecké styly (kraul, znak, prsa, motýlek)
% - Klíčové metriky: SR, SL, V = SR × SL, IdC, úhly kloubů, rotace trupu
% - Fáze záběru: catch, pull, push, recovery
% - Nejčastější chyby: dropped elbow (61.3% dle Barden & Kell 2014), špatná rotace trupu

\section{Odhad lidské pózy}
\label{sec:pose_estimation}
% TODO: Architektury: OpenPose → HRNet → ViTPose → RTMPose
% Zdroje: Xu 2022 (ViTPose), Jiang 2023 (RTMPose), Einfalt 2018/2019
% - Evoluce: bottom-up vs top-down
% - ViTPose: Vision Transformer, škálovatelný 20M-1B params
% - RTMPose: 75.8% AP COCO, 90+ FPS CPU, 430+ FPS GPU
% - Plavecky-specifické: Einfalt (Augsburg), domain-specific challenges

\section{Datasety pro plaveckou analýzu}
\label{sec:datasety}
% TODO: SwimXYZ, Augsburg, SwimmerNET, CADDY
% Zdroje: Fiche 2023 (SwimXYZ), Giulietti 2023 (SwimmerNET)
% - SwimXYZ: syntetický, 11,520 videí, 3.4M snímků, 48 keypointů
% - Augsburg: reálný, 4 styly, na žádost
% - SwimmerNET: podvodní, kraul, 2,021 snímků
% - Tabulka srovnání datasetů

\section{Strojové učení ve sportovní analýze}
\label{sec:ml_sport}
% TODO: Existující systémy, co chybí
% Diskuse RAG/LLM přístupu: Comendant 2024, Talking Tennis, BoxingPro
% Argumentace proč ML pipeline místo LLM:
% - Závislost na ext. API, nedeterminismus, latence, náklady, offline, reprodukovatelnost
% - Porovnání trade-offs: pravidlový systém vs. LLM-generated feedback

\section{Podvodní výzvy}
\label{sec:podvodni_vyzvy}
% TODO: Refrakce, bubliny, viditelnost
% - Optické artefakty pod vodou
% - Problémy s detekcí pózy (okluze, reflexy)
% - Existující přístupy: SwimmerNET, DeepLabCut
